\documentclass[11pt]{article}
\usepackage[T1]{fontenc}
\usepackage[utf8]{inputenc}
\usepackage[margin=1in]{geometry}
\usepackage{amsmath,amssymb}
\usepackage[hidelinks]{hyperref}
\title{Research scratchpad}
\author{}
\date{}
\begin{document}
\maketitle

This is a working memory for useful observations, ideas, questions and reminders.
Entries may be unverified. Supported conclusions belong in the research record,
\texttt{main.tex}; execution history belongs in \texttt{research-log.md}.

\section*{Saved entries}
No entries have been recorded yet.

% Replace the empty-state sentence when adding the first entry.
% Keep existing IDs and user notes. Example structure, not a research claim:
% \subsection*{SP-001: Short memorable title}
% \label{sp:001}
% \textbf{Recorded:} actual date. \quad \textbf{Status:} idea / observed /
% unresolved / superseded / promoted.
% Write the insight and why it is useful. Label inference and uncertainty.
% \textbf{Evidence or context:} source URL plus locator, file path,
% mind-map topic ID or experiment attempt ID when relevant.
% \textbf{Revisit when:} the question, branch or decision that needs this note.
% \textbf{Follow-up:} dated resolution or destination in the main record.

\end{document}
